% TOP_PROBLEM: {"id": 3248, "title": "Bounding the on-line choice number in terms of the choice number", "category_id": 3, "category": {"id": 3, "name": "graph_theory", "display_name": "Graph Theory", "description": "Problems involving graphs, networks, and their properties.", "slug": "graph-theory", "order_index": 3, "created_at": "2026-07-31T15:26:25.671Z"}, "status": "open", "problem_number": "OPG-48583", "set_id": 12}
% FIRST_POSTED: 2026-10-02
% ATTEMPT_STATUS: unresolved
% UnsolvedMath ID 3248; source catalogue code OPG-48583
% !TeX program = lualatex
% Research attempt by GPT-6 Astra Ultra; no claim of a new complete solution.
\documentclass[11pt,a4paper]{article}
\usepackage[margin=25mm]{geometry}
\usepackage{fontspec}
\setmainfont{lmroman10-regular.otf}[BoldFont=lmroman10-bold.otf,
 ItalicFont=lmroman10-italic.otf,BoldItalicFont=lmroman10-bolditalic.otf]
\usepackage{amsmath,amssymb,mathtools}
\usepackage{unicode-math}
\setmathfont{latinmodern-math.otf}
\AtBeginDocument{\let\setminus\smallsetminus}
\usepackage{xurl,hyperref}
\hypersetup{colorlinks=true,urlcolor=blue}
\setcounter{secnumdepth}{0}
\setlength{\parindent}{0pt}
\setlength{\parskip}{5pt}
\setlength{\emergencystretch}{3em}
\begin{document}
\section*{Bounding the on-line choice number in terms of the choice number}\label{3248}
{\small UnsolvedMath ID 3248 · OPG-48583 · Graph Theory · OpenGarden}
\subsection{Definitions and mathematical statement}
For a finite simple graph $G$, its choice number $\operatorname{ch}(G)$ is
the least $k$ such that every assignment of $k$ permitted colours per
vertex admits a proper colouring from those lists.

Define the online version through the following finite game. Initially
every vertex is uncoloured and has $k$ tokens. In each round Lister marks
a nonempty subset $M$ of the uncoloured vertices and removes one token
from each marked vertex. Painter then selects an independent subset
$I\subseteq M$ and permanently colours those vertices with the fresh
colour assigned to that round. If any vertex remains uncoloured with
zero tokens after Painter's response, Lister wins. Painter wins if all
vertices become coloured. Strategies may depend on the complete history;
neither player sees the other's future decisions.

The online choice number $\operatorname{ch}^{\mathrm{OL}}(G)$ is the least
$k$ for which Painter has a winning strategy against every play by Lister.
The independent-set requirement guarantees a proper colouring because
different rounds use different colours. This is also called the paint
number.

The source question is whether the additive gap is unbounded:
\[
 \forall b\in\mathbb N\ \exists\text{ a finite simple graph }G:
 \quad\operatorname{ch}^{\mathrm{OL}}(G)-\operatorname{ch}(G)\ge b.
\]
The graph may depend on $b$. A negative answer would give a single
absolute additive bound valid for all graphs. This question asks about
the difference, not whether the online parameter is bounded by some
arbitrary function of the offline parameter.
\subsection{Short English statement}
Can the online list-colouring requirement exceed the ordinary choice number by an arbitrarily large additive amount?
\subsection{Research attempt}
\paragraph{Scope and process.}
This is a GPT-6 Astra Ultra research attempt dated 2 October 2026, using
primary-source browsing and elementary mathematical checks. The canonical
definition above is retained. Results below are restricted results or
reductions, not a claim of a new solution to the entire catalogue question.
No formal proof assistant or external mathematical referee checked this text.


\paragraph{Literature correction and independent scope.}
The historical question in [S2] was answered affirmatively by Duraj,
Gutowski and Kozik in 2016 [R1]. Their published result gives an
unbounded additive gap already on balanced complete bipartite graphs;
it is not merely a bound for a fractional or different online parameter.
The retained catalogue ``open'' metadata is therefore historical.
This attempt proves a useful online strategy and an exact obstruction
to amplifying a finite example by disjoint unions. It does not claim
a new solution or independently reproduce the paper's difficult lower
bound for the paint number.

\paragraph{Checking the implication from the published estimates.}
Writing $p(N)$ and $c(N)$ for the paint and choice numbers of
$K_{N,N}$, the paper proves $p(N)=\log_2N+O(1)$ and derives
$c(N)\le\log_2N-(\tfrac12+o(1))\log_2\log_2N$.
Subtracting gives $p(N)-c(N)\to\infty$, which has exactly the
quantifier order requested: for each fixed integer $b$, some sufficiently
large $N$ yields a gap at least $b$. These are explicitly external
theorem inputs, not estimates inferred from numerical data or from the
leading equivalence $c(N)\sim\log_2N$. That leading equivalence alone
would not imply any growing difference.

\paragraph{A self-contained online upper strategy.}
More generally let $G$ be any bipartite graph on $N$ vertices, with
parts $X,Y$, and give each vertex $k$ tokens, where $N2^{-k}<1$.
For every uncoloured vertex with $t$ tokens remaining, assign weight
$2^{-t}$. Let $W$ be the sum of those weights. Initially
$W=N2^{-k}<1$.

In a round, let $U$ be the total old weight of marked vertices in
$X$, and $V$ the corresponding total in $Y$. Their weights double
when Lister removes a token. Painter colours all marked vertices in
the part with greater old weight, breaking ties arbitrarily. This is
an independent set. If $U\ge V$, the change in total weight is
\[
 (U+V)-2U=V-U\le0;
\]
the other case is symmetric. Hence after every response $W<1$.
An uncoloured vertex with zero tokens would contribute weight one,
which is impossible. Lister therefore cannot win.

The play also cannot continue forever without colouring all vertices.
Each nonterminal round removes at least one token, and surviving
uncoloured vertices cannot exhaust their tokens. There are finitely
many initial tokens. Thus the strategy is winning and proves
\[
 \operatorname{ch}^{\mathrm{OL}}(G)
 \le\lfloor\log_2N\rfloor+1
 \quad\text{for every nonempty bipartite }G.
\]
For $K_{N,N}$, replacing the total order by $2N$ gives an upper
bound $\log_2N+O(1)$, agreeing with the upper side of the cited
asymptotic theorem. It supplies no corresponding lower bound against
all possible Painter strategies.

\paragraph{Why separate copies cannot amplify an additive gap.}
If $G$ is the disjoint union of graphs $G_1,\ldots,G_s$, then
\[
 \operatorname{ch}(G)=\max_i\operatorname{ch}(G_i),\qquad
 \operatorname{ch}^{\mathrm{OL}}(G)=
 \max_i\operatorname{ch}^{\mathrm{OL}}(G_i).
\]
For the offline statement, colour each component from its lists
separately; the reverse inequality follows by restriction. Online,
run a winning strategy independently in each component whenever that
component receives at least one marked vertex. The union of their
responses is independent, and each component sees a legal subsequence
of rounds with exactly its own token decrements. Conversely Lister
may mark vertices only in one chosen component, reproducing any winning
play there. This proves both equalities without assuming that different
components must be assigned disjoint palettes.

In particular, taking many copies of a graph with gap one preserves
its two parameters and keeps gap one. A proof of unboundedness must
change the internal structure or use a family whose quantitative
estimates separate; increasing the number of independent copies does
not suffice.

\paragraph{Verification and final status.}
The weight invariant is checked after Painter's response, matching the
zero-token rule in the canonical definition. All marked vertices on
the selected side may be coloured together even when the graph is
not complete bipartite. The strict initial inequality prevents equality
at weight one. \textbf{Independent full derivation: UNRESOLVED.}
The precise missing component here is the growing lower separation
between the online and offline parameters. The catalogue question
itself is already answered affirmatively by [R1], whose result and
attribution must not be replaced by an open-status claim.

\subsection{Sources}
\begin{itemize}
\item[S1] ULAM.AI and the UnsolvedMath Contributors, supplied problems.json, record 3248 (OPG-48583), OpenGarden collection. Catalogue identity and source transcription; CC BY 4.0.
\url{https://huggingface.co/datasets/ulamai/UnsolvedMath}.
\item[S2] Open Problem Garden, Bounding the on-line choice number in terms of the choice number. Original problem and discussion; the mathematical formulation below is expanded from the supplied source record.
\url{http://www.openproblemgarden.org/op/bounding_the_on_line_choice_number_in_terms_of_the_choice_number}.

\item[R1] Lech Duraj, Grzegorz Gutowski and Jakub Kozik,
\emph{Chip games and paintability}, Electronic Journal of Combinatorics
23(3) (2016), P3.3. Published primary paper, Theorem 1 and Corollary 2
checked 2 October 2026.
\url{https://www.combinatorics.org/ojs/index.php/eljc/article/download/v23i3p3/pdf/}.

\end{itemize}
\end{document}
